\section{张量并行（Tensor Parallelism）}

与流水线并行"纵向"切分不同，张量并行是"横向"切分：\textbf{将每一层的矩阵乘法分解为子矩阵，分配到不同 GPU 上并行计算}。

\subsection{基本原理：矩阵乘法分块}

\begin{figure}[H]
\centering
\includegraphics[width=0.95\textwidth]{slides-images/slide-038.jpg}
\caption{张量并行的数学基础：将 $Y = XA$ 分解为 $Y = X_1 A_1 + X_2 A_2$\protect\footnotemark}
\end{figure}
\footnotetext{Slides 第39页。}

对于矩阵乘法 $Y = XA$，可以沿不同维度切分矩阵 $A$：

$$Y = X \cdot A = X \cdot [A_1 | A_2] = [XA_1 | XA_2]  \quad \text{（列切分）}$$

$$Y = X \cdot A = [X_1 | X_2] \cdot \begin{bmatrix} A_1 \\ A_2 \end{bmatrix} = X_1 A_1 + X_2 A_2  \quad \text{（行切分）}$$

每个 GPU 负责一个子矩阵的乘法，然后通过 All-Reduce 或拼接来聚合结果。

\subsection{MLP 中的张量并行}

\begin{figure}[H]
\centering
\includegraphics[width=0.95\textwidth]{slides-images/slide-039.jpg}
\caption{MLP 的张量并行实现：f 为复制/All-Reduce 同步点，g 为 All-Reduce/复制同步点\protect\footnotemark}
\end{figure}
\footnotetext{Slides 第40页。}

以两层 MLP（$Z = \text{dropout}(\text{GeLU}(XA) \cdot B)$）为例：

\begin{enumerate}
    \item \textbf{前向传播}：
    \begin{itemize}
        \item 在 $f$ 处复制输入 $X$ 到所有 GPU
        \item 每个 GPU 计算 $Y_i = \text{GeLU}(X \cdot A_i)$，然后 $Z_i = Y_i \cdot B_i$
        \item 在 $g$ 处 All-Reduce 求和：$Z = \sum_i Z_i$
    \end{itemize}
    \item \textbf{反向传播}：
    \begin{itemize}
        \item 在 $g$ 处复制梯度到所有 GPU
        \item 各 GPU 独立计算反向
        \item 在 $f$ 处 All-Reduce 梯度
    \end{itemize}
\end{enumerate}

\begin{importantbox}{张量并行的同步模式}
\textbf{每一层}需要一次前向 All-Reduce 和一次反向 All-Reduce。对于 $L$ 层 Transformer，每个训练步需要 $2L$ 次 All-Reduce，每次通信量为 $B \times S \times H$（批次 $\times$ 序列长度 $\times$ 隐藏维度）。这是一种通信密集型策略。
\end{importantbox}

\subsection{张量并行的经验法则}

\begin{figure}[H]
\centering
\includegraphics[width=0.95\textwidth]{slides-images/slide-040.jpg}
\caption{不同张量并行度下的吞吐量变化\protect\footnotemark}
\end{figure}
\footnotetext{Slides 第41页。来自 HuggingFace 并行化教程。}

\begin{warningbox}{张量并行度 $\leq$ 8（节点内 GPU 数）}
实验数据显示：
\begin{itemize}
    \item TP=2: 吞吐下降约 10\%
    \item TP=4: 吞吐下降约 10\%
    \item TP=8: 吞吐下降约 12\%（仍在节点内）
    \item TP=16: 吞吐下降约 \textbf{42\%}（跨节点！）
    \item TP=32: 吞吐下降约 \textbf{65\%}
\end{itemize}
一旦超出节点内的 NVLink 带宽，性能急剧恶化。因此，\textbf{张量并行应始终限制在单节点内}。
\end{warningbox}

\subsection{张量并行 vs. 流水线并行}

\begin{figure}[H]
\centering
\includegraphics[width=0.95\textwidth]{slides-images/slide-042.jpg}
\caption{张量并行与流水线并行的对比分析\protect\footnotemark}
\end{figure}
\footnotetext{Slides 第43页。}

\begin{table}[H]
\centering
\begin{tabular}{lcc}
\toprule
\textbf{特性} & \textbf{流水线并行} & \textbf{张量并行} \\
\midrule
切分方式 & 按层（深度） & 按矩阵（宽度） \\
通信类型 & 点对点 & All-Reduce \\
通信量/层 & $B \times S \times H$ & $8 \times B \times S \times H$ \\
气泡 & 有（需微批次缓解） & 无 \\
批次消耗 & 是 & 否 \\
实现复杂度 & 非常高 & 中等 \\
最佳使用场景 & 跨节点慢速链路 & 节点内高速互连 \\
\bottomrule
\end{tabular}
\caption{两种模型并行策略的关键区别}
\end{table}

\subsection{本章小结}

张量并行通过矩阵分块将单层计算分布到多个 GPU 上，无气泡开销但通信密集。实践中应限制在单节点内（$\leq 8$ GPU），利用 NVLink 高带宽。与流水线并行互补：张量并行用于节点内，流水线并行用于跨节点。

%% ========== Part 6: 序列并行与激活分片 ==========

